\subsection{The Floating-Point Representation Architecture (\texttt{float})}

Python's \texttt{float} type represents real-number approximations, not exact real numbers. This distinction is essential. An \texttt{int} can represent arbitrarily large whole numbers exactly. A \texttt{float} represents numbers using a fixed-size binary format, so many decimal fractions can only be approximated.

A first inspection:

\begin{verbatim}
price = 19.99

print(f"value: {price}")
print(f"type:  {type(price)}")
print(f"id:    {id(price)}")

print()
print("selected float methods:")
print(f"{'as_integer_ratio':<20} {'as_integer_ratio' in dir(price)}")
print(f"{'is_integer':<20} {'is_integer' in dir(price)}")
print(f"{'hex':<20} {'hex' in dir(price)}")
print(f"{'__add__':<20} {'__add__' in dir(price)}")
print(f"{'__truediv__':<20} {'__truediv__' in dir(price)}")

print()
print(f"is_integer():       {price.is_integer()}")
print(f"as_integer_ratio(): {price.as_integer_ratio()}")
print(f"hex():              {price.hex()}")
\end{verbatim}

Possible output:

\begin{verbatim}
value: 19.99
type:  <class 'float'>
id:    2214135699408

selected float methods:
as_integer_ratio     True
is_integer           True
hex                  True
__add__              True
__truediv__          True

is_integer():       False
as_integer_ratio(): (5626684784442573, 281474976710656)
hex():              0x1.3fd70a3d70a3dp+4
\end{verbatim}

The method \texttt{as\_integer\_ratio()} exposes the exact fraction represented by the stored floating-point value. The method \texttt{hex()} exposes a hexadecimal form of the binary floating-point value. These are useful because they show that \texttt{19.99} is not stored as the exact decimal fraction that the programmer wrote.

\subsubsection{Architectural Realities: Mapping Python Floats to C Doubles via the IEEE 754 Double-Precision Binary Standard}

In CPython, the ordinary Python \texttt{float} is implemented using the platform C \texttt{double}. On modern machines, that usually means IEEE 754 double-precision binary floating point.

This format uses a fixed amount of storage. It does not grow like Python's \texttt{int}. A float has limited precision and limited exponent range.

Python exposes the main limits through \texttt{sys.float\_info}:

\begin{verbatim}
import sys

print(sys.float_info)

print()
print(f"max:      {sys.float_info.max}")
print(f"min:      {sys.float_info.min}")
print(f"epsilon:  {sys.float_info.epsilon}")
print(f"mantissa digits in base 2: {sys.float_info.mant_dig}")
\end{verbatim}

Possible output:

\begin{verbatim}
sys.float_info(max=1.7976931348623157e+308, min=2.2250738585072014e-308, 
dig=15, mant_dig=53, epsilon=2.220446049250313e-16, radix=2, rounds=1)

max:      1.7976931348623157e+308
min:      2.2250738585072014e-308
epsilon:  2.220446049250313e-16
mantissa digits in base 2: 53
\end{verbatim}

The exact formatting may vary, but the important values are stable on normal CPython builds:

\begin{itemize}
    \item \texttt{radix = 2}: the representation is binary;
    \item \texttt{mant\_dig = 53}: the precision is about 53 binary digits;
    \item \texttt{dig = 15}: about 15 decimal digits can be represented reliably;
    \item \texttt{epsilon}: the distance between \texttt{1.0} and the next representable float above it.
\end{itemize}

This means that \texttt{float} is not an arbitrary-precision decimal type. It is a fixed-precision binary approximation type.

\begin{verbatim}
small = 42.0
large = 10.0 ** 100

print(f"small: {small}")
print(f"large: {large}")

print()
print(f"type(small): {type(small)}")
print(f"type(large): {type(large)}")

print()
print(f"small.hex(): {small.hex()}")
print(f"large.hex(): {large.hex()}")
\end{verbatim}

Possible output:

\begin{verbatim}
small: 42.0
large: 1e+100

type(small): <class 'float'>
type(large): <class 'float'>

small.hex(): 0x1.5000000000000p+5
large.hex(): 0x1.249ad2594c37dp+332
\end{verbatim}

Both values have type \texttt{float}. The second value is much larger, but the representation is still constrained by the same fixed floating-point architecture.

This also means that float arithmetic can overflow:

\begin{verbatim}
import sys

big = sys.float_info.max

print(f"big: {big}")

too_big = big * 2

print(f"too_big: {too_big}")
print(f"type(too_big): {type(too_big)}")
\end{verbatim}

Possible output:

\begin{verbatim}
big: 1.7976931348623157e+308
too_big: inf
type(too_big): <class 'float'>
\end{verbatim}

Unlike \texttt{int}, a \texttt{float} cannot keep allocating more storage to represent a larger exact value. Once the representable range is exceeded, the result may become infinity.

\subsubsection{Machine Epsilon and Precision Loss: The Pitfalls of Representing Base-10 Fractions in Base-2 Memory Buffers}

The most famous floating-point surprise is this:

\begin{verbatim}
total = 0.1 + 0.2

print(f"total: {total}")
print(f"total == 0.3: {total == 0.3}")
\end{verbatim}

Possible output:

\begin{verbatim}
total: 0.30000000000000004
total == 0.3: False
\end{verbatim}

This is not a Python bug. It is a consequence of representing decimal fractions in binary floating point.

Fractions such as \texttt{0.5}, \texttt{0.25}, and \texttt{0.125} are easy in binary because they are powers of two in the denominator:

\begin{verbatim}
values = [0.5, 0.25, 0.125]

for value in values:
    print(f"{value:<8} ratio={value.as_integer_ratio()}")
\end{verbatim}

Possible output:

\begin{verbatim}
0.5      ratio=(1, 2)
0.25     ratio=(1, 4)
0.125    ratio=(1, 8)
\end{verbatim}

But \texttt{0.1} does not have a finite binary expansion, just as \texttt{1/3} does not have a finite decimal expansion.

\begin{verbatim}
value = 0.1

print(f"value: {value}")
print(f"ratio: {value.as_integer_ratio()}")
print(f"hex:   {value.hex()}")
\end{verbatim}

Possible output:

\begin{verbatim}
value: 0.1
ratio: (3602879701896397, 36028797018963968)
hex:   0x1.999999999999ap-4
\end{verbatim}

The source text \texttt{0.1} is converted to the closest representable binary floating-point value. That stored value is very close to one tenth, but it is not exactly one tenth.

Machine epsilon gives a scale for this limitation around \texttt{1.0}.

\begin{verbatim}
import sys

eps = sys.float_info.epsilon

print(f"epsilon: {eps}")
print(f"1.0 + epsilon:     {1.0 + eps}")
print(f"1.0 + epsilon / 2: {1.0 + eps / 2}")

print()
print(f"(1.0 + epsilon) == 1.0:     {(1.0 + eps) == 1.0}")
print(f"(1.0 + epsilon / 2) == 1.0: {(1.0 + eps / 2) == 1.0}")
\end{verbatim}

Possible output:

\begin{verbatim}
epsilon: 2.220446049250313e-16
1.0 + epsilon:     1.0000000000000002
1.0 + epsilon / 2: 1.0

(1.0 + epsilon) == 1.0:     False
(1.0 + epsilon / 2) == 1.0: True
\end{verbatim}

The second addition is too small to change the stored value at that scale. This is the concrete meaning of limited precision: not every tiny mathematical difference can be represented.

For equality checks with floats, direct equality is often wrong for computed values:

\begin{verbatim}
total = 0.1 + 0.2
expected = 0.3

print(f"total:    {total}")
print(f"expected: {expected}")
print(f"direct equality: {total == expected}")

print()
print(f"difference: {abs(total - expected)}")
\end{verbatim}

Possible output:

\begin{verbatim}
total:    0.30000000000000004
expected: 0.3
direct equality: False

difference: 5.551115123125783e-17
\end{verbatim}

A better test is to check whether the values are close enough for the problem being solved:

\begin{verbatim}
import math

total = 0.1 + 0.2
expected = 0.3

print(f"math.isclose(total, expected): {math.isclose(total, expected)}")
\end{verbatim}

Possible output:

\begin{verbatim}
math.isclose(total, expected): True
\end{verbatim}

The practical rule is:

\begin{quote}
Use \texttt{float} for approximate real-number computation, not for exact decimal accounting.
\end{quote}

For money, exact decimal rules, or base-10 financial calculations, Python's standard \texttt{decimal} module is often more appropriate, but that belongs to a later discussion.

\subsubsection{Special Floating-Point Values: Infinity, Negative Infinity, and NaN}

Floating-point arithmetic also includes special values. The most important are:

\begin{itemize}
    \item positive infinity: \texttt{inf};
    \item negative infinity: \texttt{-inf};
    \item not-a-number: \texttt{nan}.
\end{itemize}

They can be created with \texttt{float()} or accessed through the \texttt{math} module.

\begin{verbatim}
import math

pos_inf = float("inf")
neg_inf = float("-inf")
not_number = float("nan")

print(f"pos_inf:    {pos_inf}")
print(f"neg_inf:    {neg_inf}")
print(f"not_number: {not_number}")

print()
print(f"type(pos_inf):    {type(pos_inf)}")
print(f"type(neg_inf):    {type(neg_inf)}")
print(f"type(not_number): {type(not_number)}")
\end{verbatim}

Possible output:

\begin{verbatim}
pos_inf:    inf
neg_inf:    -inf
not_number: nan

type(pos_inf):    <class 'float'>
type(neg_inf):    <class 'float'>
type(not_number): <class 'float'>
\end{verbatim}

These are still \texttt{float} objects. They are not separate Python types.

Infinity compares larger than ordinary finite numbers:

\begin{verbatim}
pos_inf = float("inf")
neg_inf = float("-inf")

print(f"pos_inf > 10 ** 300: {pos_inf > 10 ** 300}")
print(f"neg_inf < -10 ** 300: {neg_inf < -10 ** 300}")

print()
print(f"pos_inf + 42: {pos_inf + 42}")
print(f"neg_inf * 2:  {neg_inf * 2}")
\end{verbatim}

Possible output:

\begin{verbatim}
pos_inf > 10 ** 300: True
neg_inf < -10 ** 300: True

pos_inf + 42: inf
neg_inf * 2:  -inf
\end{verbatim}

The value \texttt{nan} is different. It represents an undefined or unrepresentable floating-point result.

\begin{verbatim}
not_number = float("nan")

print(f"not_number: {not_number}")
print(f"not_number == not_number: {not_number == not_number}")
print(f"not_number != not_number: {not_number != not_number}")
\end{verbatim}

Possible output:

\begin{verbatim}
not_number: nan
not_number == not_number: False
not_number != not_number: True
\end{verbatim}

This is strange but intentional. A NaN is not considered equal even to itself. Therefore, do not test for NaN using equality.

Use \texttt{math.isnan()}:

\begin{verbatim}
import math

value1 = 42.0
value2 = float("nan")

print(f"value1 is NaN: {math.isnan(value1)}")
print(f"value2 is NaN: {math.isnan(value2)}")
\end{verbatim}

Possible output:

\begin{verbatim}
value1 is NaN: False
value2 is NaN: True
\end{verbatim}

Similarly, use \texttt{math.isinf()} to test for infinity:

\begin{verbatim}
import math

values = [42.0, float("inf"), float("-inf"), float("nan")]

for value in values:
    print(f"{str(value):<8} "
          f"isinf={math.isinf(value)!s:<5} "
          f"isnan={math.isnan(value)!s:<5} "
          f"isfinite={math.isfinite(value)!s:<5}")
\end{verbatim}

Possible output:

\begin{verbatim}
42.0     isinf=False isnan=False isfinite=True 
inf      isinf=True  isnan=False isfinite=False
-inf     isinf=True  isnan=False isfinite=False
nan      isinf=False isnan=True  isfinite=False
\end{verbatim}

The practical summary is:

\begin{quote}
Python's \texttt{float} is a fixed-precision binary floating-point object. It is fast and useful for approximate numerical computation, but it cannot represent all decimal fractions exactly, can overflow to infinity, and includes special values such as \texttt{inf}, \texttt{-inf}, and \texttt{nan}.
\end{quote}